% Introduction. First draft, Diego with Claude, 22 Nov, on the frozen numbers. Diego and Martin write the
% Introduction and Discussion together (DECISIONS G3): Martin to revise. Budget about 0.75 page.

\dropcap{T}he United States Congress is usually described as two camps, and most ways of measuring it agree.
Members of the two parties vote apart, and they have used increasingly different words since the early 1990s
\cite{gentzkow2019}. Networks built from who signs whose bills show the same split, with ties between the parties
growing rarer and more hostile over four decades \cite{neal2020}.

How deep the split looks depends, however, on which record is read. Between 1973 and 2004, bipartisan agreement
fell in roll-call votes while staying far more stable in cosponsorship, because party leaders chose which bills
reached a vote \cite{harbridge2015}. Across communication channels the parties are always distinguishable, but more
so in front of national audiences than local ones \cite{green2024}. Division is a property of the measurement as
well as of the institution.

Members also have reasons to manage how divided they appear. Voters with weak party attachments reward a
bipartisan image \cite{harbridge2011}, and what members say to constituents is a strategic presentation of their
work \cite{grimmer2013}. That presentation need not match the work: floor speeches that call a bill bipartisan
mostly concern bills with trivial support from the other party, and a member's bipartisan rhetoric is unrelated to
their bipartisan votes \cite{westwood2022}.

We compare two records of the same 450 members of the House in the 119th Congress (2025 to 2026). The first is the
cosponsorship network \cite{fowler2006}: two members are tied when one signs a bill the other sponsors. The second
is, to our knowledge, new: two members are tied when one names the other in a floor speech. Both are networks on
identical nodes, so the same measures of party division apply to each and both can be tested against the same
null models. A word such as ``bipartisan'' costs nothing to say. A name points at a specific colleague, and what is
said about that colleague can be read.

Whether naming is more or less partisan than working cannot be guessed in advance. Members might attack the other
side in public and cooperate quietly, or display the few cross-party partners they have. We find that the floor map
is less divided by party than the working map at equal size (38.7\% of ties cross the aisle, against 28.6\%), that
debate procedure accounts for much, but not all, of the floor map's own party division, and that cross-party mentions are mostly
acknowledgement rather than argument. At home, in press releases, members in exposed seats name the other party
more often than members in safe seats.
